\chapter{Zwei Irrfahrten, ein Band}

\section{Der Ausnahmefall}

\begin{merksatz}
\boxtitel{Arzneipreise und Arzneiausgaben}
Beide Reihen --- der Preisindex und die Ausgaben je Versicherten, logarithmiert, zwanzig Jahre monatlich --- haben eine Einheitswurzel. Nach Kapitel~11 müsste man differenzieren.

\begin{center}
{\small
\renewcommand{\arraystretch}{1.48}
\begin{tabular}{@{}P{3.6cm}P{2.4cm}@{}}
\toprule
\textbf{Regression} & \textbf{Elastizität}\\
\midrule
in Niveaus & $0{,}575$\\
in Differenzen & $0{,}506$\\
\midrule
wahrer Langfristwert & $\mathbf{0{,}600}$\\
\bottomrule
\end{tabular}}
\end{center}

\begin{center}
\begin{minipage}{0.9\textwidth}
\textbf{Hier ist die Niveauregression die bessere --- obwohl beide Reihen nichtstationär sind. Warum?}
\end{minipage}
\end{center}

\begin{center}
\begin{minipage}{0.88\textwidth}
\kursiv{Weil die Ausgaben den Preisen nicht beliebig davonlaufen können: Eine Zeitlang schon --- Nachfrageverschiebungen, neue Medikamente ---, aber auf Dauer bindet die Preisregulierung beide aneinander. \textbf{Die Abweichung vom Gleichgewicht ist stationär, obwohl die Reihen selbst es nicht sind.} Das ist der Ausnahmefall, der die Regel von Kapitel~11 durchbricht.}
\end{minipage}
\end{center}
\end{merksatz}

\begin{definition}
\boxtitelm{Definition 12.1 --- Kointegration}
Zwei Reihen $x_t$, $y_t$ mit je einer Einheitswurzel heißen \Begriff{kointegriert}, wenn es ein $\beta$ gibt, sodass
\[ u_t = y_t - \beta\,x_t \]
stationär ist. Der Vektor $(1,-\beta)$ heißt Kointegrationsvektor, $u_t$ der Gleichgewichtsfehler.
\end{definition}

\begin{merksatz}
\boxtitel{Satz 12.2 --- Superkonsistenz}
Sind $x_t$ und $y_t$ kointegriert, so strebt der Kleinste-Quadrate-Schätzer von $\beta$ aus der Niveauregression gegen den wahren Wert --- und zwar mit Rate $1/T$ statt $1/\sqrt T$.

\schritt{Der Gedanke}
\[ \hat\beta - \beta = \frac{\sum x_tu_t}{\sum x_t^{2}} \]
Der Nenner wächst wie $T^{2}$ (Irrfahrt, Band 25), der Zähler nur wie $T$ --- $u_t$ ist stationär und $x_t$ läuft davon, ihre Produkte summieren sich langsamer. Der Quotient fällt wie $1/T$. $\blacksquare$

\schritt{Simuliert}
\begin{center}
{\small
\renewcommand{\arraystretch}{1.48}
\begin{tabular}{@{}P{1.6cm}P{2.6cm}P{2.4cm}@{}}
\toprule
$T$ & \textbf{mittl.\ Fehler} & \textbf{Fehler}$\,\cdot\,T$\\
\midrule
$50$ & $0{,}0218$ & $1{,}09$\\
$100$ & $0{,}0115$ & $1{,}15$\\
$400$ & $0{,}0031$ & $1{,}24$\\
$1600$ & $0{,}0007$ & $1{,}20$\\
\bottomrule
\end{tabular}}
\end{center}

\begin{center}
\begin{minipage}{0.9\textwidth}
\textbf{Vervierfacht man die Daten, viertelt sich der Fehler --- nicht halbiert. Die Nichtstationarität hilft, statt zu schaden.}
\end{minipage}
\end{center}

\begin{center}
\begin{minipage}{0.88\textwidth}
\kursiv{\textbf{Die Endogenität von $x$ spielt dabei keine Rolle}: Selbst wenn $x_t$ mit $u_t$ korreliert ist, bleibt der Schätzer konsistent, weil der Nenner alles überrollt. Das ist einer der wenigen Fälle in diesem Band, in denen die kleinsten Quadrate trotz Endogenität richtig liegen --- und der einzige, in dem der Grund die Nichtstationarität ist. Die Standardfehler sind allerdings unbrauchbar; dafür gibt es eigene Verfahren.}
\end{minipage}
\end{center}
\end{merksatz}

\section{Kointegration prüfen}

\begin{merksatz}
\boxtitel{Satz 12.3 --- der Test von Engle und Granger}
\begin{center}
{\small
\renewcommand{\arraystretch}{1.55}
\begin{tabular}{@{}P{1.0cm}P{7.4cm}@{}}
\toprule
$1$ & Niveauregression $y$ auf $x$; Residuen $\hat u_t$ aufheben.\\
$2$ & Dickey-Fuller-Test auf $\hat u_t$ --- mit \kursiv{eigenen} kritischen Werten, weil $\hat\beta$ geschätzt ist.\\
$3$ & Wird die Einheitswurzel der Residuen verworfen, sind die Reihen kointegriert.\\
\bottomrule
\end{tabular}}
\end{center}

\begin{center}
{\small
\renewcommand{\arraystretch}{1.48}
\begin{tabular}{@{}P{3.6cm}P{2.0cm}P{2.4cm}@{}}
\toprule
\textbf{Lage} & $t$ & \textbf{Urteil}\\
\midrule
Preise und Ausgaben & $-7{,}69$ & \textcolor{gruen}{kointegriert}\\
unabhängige Irrfahrten & verwirft in $4{,}8\,\%$ & \textcolor{gruen}{Test hält}\\
\midrule
kritischer Wert, $5\,\%$ & $-3{,}37$ & \\
\bottomrule
\end{tabular}}
\end{center}

\begin{center}
\begin{minipage}{0.88\textwidth}
\kursiv{\textbf{Der kritische Wert ist strenger als beim gewöhnlichen Dickey-Fuller-Test ($-2{,}89$)}, weil die kleinsten Quadrate die Residuen so stationär wie möglich machen --- auch dann, wenn sie es nicht sind. Ohne die Berichtigung fände man bei unabhängigen Irrfahrten in etwa einem Fünftel der Fälle Kointegration.}
\end{minipage}
\end{center}
\end{merksatz}

\section{Kurz- und Langfrist}

\begin{merksatz}
\boxtitel{Satz 12.4 --- das Fehlerkorrekturmodell}
Sind $x_t$ und $y_t$ kointegriert, so lässt sich die Dynamik schreiben als
\[ \Delta y_t = \alpha\,\underbrace{(y_{t-1}-\beta x_{t-1})}_{u_{t-1}} + \gamma\,\Delta x_t + \varepsilon_t, \qquad \alpha < 0 . \]

\begin{center}
{\small
\renewcommand{\arraystretch}{1.5}
\begin{tabular}{@{}P{1.4cm}P{7.0cm}@{}}
\toprule
$\gamma$ & kurzfristige Reaktion auf eine Preisänderung.\\
$\beta$ & langfristige Elastizität.\\
$\alpha$ & Anpassungsgeschwindigkeit: welcher Anteil der Lücke je Periode geschlossen wird.\\
\bottomrule
\end{tabular}}
\end{center}

\schritt{Der Gedanke}
Ist $y$ über dem Gleichgewicht ($u_{t-1}>0$), so muss es fallen --- sonst wäre $u$ nicht stationär. Der Term $\alpha u_{t-1}$ mit $\alpha<0$ erzwingt die Rückkehr. Dass Kointegration und Fehlerkorrektur dasselbe sind, ist der Darstellungssatz von Granger; wir übergehen den Beweis --- die dritte Auslassung dieses Bandes. $\blacksquare$

\schritt{Im Beispiel}
\begin{center}
{\small
\renewcommand{\arraystretch}{1.48}
\begin{tabular}{@{}P{3.0cm}P{2.0cm}P{3.0cm}@{}}
\toprule
\textbf{Parameter} & \textbf{Schätzung} & \textbf{Deutung}\\
\midrule
$\gamma$ (kurz) & $0{,}52$ & sofort gut die Hälfte\\
$\beta$ (lang) & $0{,}58$ & auf Dauer $0{,}6$\\
$\alpha$ (Anpassung) & $\mathbf{-0{,}39}$ & $39\,\%$ der Lücke je Monat\\
\bottomrule
\end{tabular}}
\end{center}

\schritt{Ein Preissprung von zehn Prozent}
\begin{center}
{\small
\renewcommand{\arraystretch}{1.46}
\begin{tabular}{@{}P{2.0cm}P{1.4cm}P{1.4cm}P{1.4cm}P{1.4cm}P{1.4cm}@{}}
\toprule
\textbf{Monate} & $0$ & $1$ & $2$ & $4$ & $8$\\
\midrule
offene Lücke & $6{,}0\,\%$ & $3{,}7$ & $2{,}2$ & $0{,}8$ & $0{,}1$\\
\bottomrule
\end{tabular}}
\end{center}

\begin{center}
\begin{minipage}{0.9\textwidth}
\textbf{Halbwertszeit der Anpassung: $1{,}4$ Monate. Nach einem halben Jahr ist das neue Gleichgewicht erreicht.}
\end{minipage}
\end{center}

\begin{center}
\begin{minipage}{0.88\textwidth}
\kursiv{\textbf{Das Fehlerkorrekturmodell versöhnt die Kapitel~10 und 11}: Es arbeitet in Differenzen, also stationär, und behält trotzdem die Langfristbeziehung --- als zurückliegenden Gleichgewichtsfehler. Für die Politik liefert es drei getrennte Antworten: wie stark sofort, wie stark auf Dauer, und wie lange dazwischen. Die Differenzenregression aus Kapitel~11 kennt nur die erste.}
\end{minipage}
\end{center}
\end{merksatz}

\section{Mehrere Reihen zugleich}

\begin{merksatz}
\boxtitel{Die Vektorautoregression}
\[ \begin{pmatrix} y_{1t}\\ y_{2t}\end{pmatrix} = A\begin{pmatrix} y_{1,t-1}\\ y_{2,t-1}\end{pmatrix} + \begin{pmatrix}\varepsilon_{1t}\\ \varepsilon_{2t}\end{pmatrix} \]
--- jede Reihe auf die Vergangenheit \kursiv{aller}, ohne Ausschlüsse.

\schritt{Granger-Kausalität}
$y_2$ \Begriff{Granger-verursacht} $y_1$, wenn die Vergangenheit von $y_2$ die Vorhersage von $y_1$ verbessert --- also der Koeffizient von $y_{2,t-1}$ in der ersten Gleichung nicht null ist.

\begin{center}
{\small
\renewcommand{\arraystretch}{1.48}
\begin{tabular}{@{}P{2.4cm}P{2.6cm}P{2.6cm}@{}}
\toprule
\textbf{Gleichung} & \textbf{auf } $y_{1,t-1}$ & \textbf{auf } $y_{2,t-1}$\\
\midrule
$y_1$ & $0{,}53$ ($t=13{,}6$) & $\mathbf{0{,}31}$ ($t=7{,}5$)\\
$y_2$ & $0{,}03$ ($t=0{,}9$) & $0{,}59$ ($t=14{,}1$)\\
\bottomrule
\end{tabular}}
\end{center}
\kursiv{Simuliert mit $A = \bigl(\begin{smallmatrix}0{,}5&0{,}3\\0&0{,}6\end{smallmatrix}\bigr)$: $y_2$ wirkt auf $y_1$, nicht umgekehrt --- und die Schätzung findet es.}

\begin{center}
\begin{minipage}{0.88\textwidth}
\kursiv{\textbf{Granger-Kausalität ist Vorhersagbarkeit, keine Kausalität im Sinne von Band 23.} Wer den Wetterbericht hört, bevor er den Schirm nimmt: Der Bericht Granger-verursacht den Schirm, und beide werden vom Wetter verursacht. Der Name ist unglücklich, und Granger hat es selbst bedauert. Die Impulsantwort --- wie $y_1$ auf einen Schock in $y_2$ reagiert --- ist dagegen eine kausale Aussage und braucht deshalb eine Identifikationsannahme: Kapitel~2 kehrt zurück, diesmal als Frage, welcher Schock zuerst wirkt.}
\end{minipage}
\end{center}
\end{merksatz}

\section{Der Bruchfall}

\begin{warnung}
\boxtitelr{Woran Kointegration scheitert}
\begin{center}
{\small
\renewcommand{\arraystretch}{1.52}
\begin{tabular}{@{}P{2.8cm}P{5.8cm}@{}}
\toprule
\textbf{Mehr als zwei Reihen} & Es kann mehrere Kointegrationsvektoren geben; Engle-Granger findet höchstens einen und hängt davon ab, welche Reihe links steht.\\
\textbf{Strukturbruch} & Ändert sich $\beta$ --- eine Reform der Preisregulierung ---, so ist der Gleichgewichtsfehler nicht mehr stationär, und der Test verwirft die Kointegration, obwohl sie abschnittsweise gilt.\\
\textbf{Nahe eins} & Ist $x$ nur fast eine Einheitswurzel, so gilt weder die Theorie von Kapitel~11 noch die von Kapitel~12 genau.\\
\bottomrule
\end{tabular}}
\end{center}

\begin{center}
\begin{minipage}{0.88\textwidth}
\kursiv{\textbf{Für mehrere Reihen gibt es das Verfahren von Johansen}, das den Rang der Langfristmatrix schätzt und damit die Zahl der Kointegrationsbeziehungen --- eine Eigenwertaufgabe aus Band 2, und eine der eleganteren Konstruktionen der Ökonometrie. Die dritte Zeile ist der unangenehmste Fall: \textbf{Zwischen $\phi = 0{,}95$ und $\phi = 1$ liegt ein Niemandsland}, in dem beide Theorien nur näherungsweise gelten und die Vertrauensintervalle von keiner stimmen.}
\end{minipage}
\end{center}
\end{warnung}

\section{Übungsaufgaben}

\begin{aufgabe}
\boxtitel{Aufgabe 12.1 --- Den Gleichgewichtsfehler deuten}
Im Beispiel seien Preise um $4$ Prozent gestiegen, die Ausgaben um $1$ Prozent. Wo steht das System, und was folgt für den nächsten Monat?

\loesung
\schritt{Schritt 1 --- Der Gleichgewichtsfehler}
\[ u = \Delta y - \beta\,\Delta x = 1 - 0{,}6\cdot4 = -1{,}4 \ \text{Prozent} \]

\schritt{Schritt 2 --- Die Korrektur}
\[ \alpha\cdot u = -0{,}39\cdot(-1{,}4) = +0{,}55 \ \text{Prozent} \]

\begin{center}
\begin{minipage}{0.9\textwidth}
\textbf{Die Ausgaben liegen $1{,}4$ Prozent unter dem Gleichgewicht und werden im nächsten Monat um etwa einen halben Prozentpunkt zusätzlich steigen --- auch ohne weitere Preisänderung.}
\end{minipage}
\end{center}

\schritt{Schritt 3 --- Die Vorhersage über mehrere Monate}
\begin{center}
{\small
\renewcommand{\arraystretch}{1.48}
\begin{tabular}{@{}P{1.6cm}P{2.6cm}P{2.6cm}@{}}
\toprule
\textbf{Monat} & \textbf{Lücke} & \textbf{Zusatzanstieg}\\
\midrule
$1$ & $-1{,}40$ & $+0{,}55$\\
$2$ & $-0{,}85$ & $+0{,}33$\\
$3$ & $-0{,}52$ & $+0{,}20$\\
$6$ & $-0{,}12$ & $+0{,}05$\\
\bottomrule
\end{tabular}}
\end{center}

\begin{center}
\begin{minipage}{0.88\textwidth}
\kursiv{\textbf{Das ist der Vorhersagewert der Kointegration}: Der Gleichgewichtsfehler heute sagt die Bewegung morgen voraus. Die Differenzenregression aus Kapitel~11 hätte für den nächsten Monat null vorhergesagt --- keine Preisänderung, keine Bewegung. Für Kassen, die Budgets planen, ist der Unterschied erheblich.}
\end{minipage}
\end{center}
\end{aufgabe}

\begin{aufgabe}
\boxtitel{Aufgabe 12.2 --- Welche Reihe steht links?}
Bei Engle-Granger hängt das Ergebnis davon ab, ob man $y$ auf $x$ oder $x$ auf $y$ regressiert. Warum, und was tut man?

\loesung
\schritt{Schritt 1 --- Der Grund}
Die kleinsten Quadrate minimieren senkrechte Abstände in Richtung der linken Variablen. Bei endlichem $T$ sind $\hat\beta_{y|x}$ und $1/\hat\beta_{x|y}$ verschieden --- die Dämpfung aus Kapitel~1, denn der Regressor ist mit Rauschen behaftet.

\schritt{Schritt 2 --- Im Beispiel}
\begin{center}
{\small
\renewcommand{\arraystretch}{1.48}
\begin{tabular}{@{}P{3.0cm}P{2.4cm}@{}}
\toprule
\textbf{Regression} & $\hat\beta$\\
\midrule
Ausgaben auf Preise & $0{,}575$\\
Preise auf Ausgaben, invertiert & $\approx 0{,}63$\\
\bottomrule
\end{tabular}}
\end{center}

\schritt{Schritt 3 --- Die Abhilfe}
\begin{center}
{\small
\renewcommand{\arraystretch}{1.5}
\begin{tabular}{@{}P{2.8cm}P{5.8cm}@{}}
\toprule
\textbf{Asymptotisch} & Satz 12.2 gilt in beide Richtungen; der Unterschied verschwindet mit $1/T$.\\
\textbf{Endlich} & Johansens Verfahren behandelt alle Reihen gleich und kennt kein Links und Rechts.\\
\textbf{Inhaltlich} & Meist ist klar, welche Größe sich anpasst --- und die steht links.\\
\bottomrule
\end{tabular}}
\end{center}

\begin{center}
\begin{minipage}{0.88\textwidth}
\kursiv{\textbf{Die dritte Zeile verbirgt eine kausale Annahme}: Dass die Ausgaben den Preisen folgen und nicht umgekehrt. Bei regulierten Preisen ist das plausibel; bei Marktpreisen ist es die Simultanität aus Kapitel~2. Das Fehlerkorrekturmodell kann prüfen, wer sich anpasst --- man schätzt $\alpha$ in beiden Gleichungen und sieht, welches signifikant negativ ist.}
\end{minipage}
\end{center}
\end{aufgabe}

\begin{aufgabe}
\boxtitel{Aufgabe 12.3 --- Impulsantworten und ihre Annahme}
Die Residuen eines VAR aus zwei Gleichungen sind mit $0{,}46$ korreliert. Was heißt das für die Frage, wie $y_1$ auf einen Schock in $y_2$ reagiert?

\loesung
\schritt{Schritt 1 --- Das Problem}
Ein Residuum ist kein Schock: Wenn $\varepsilon_1$ und $\varepsilon_2$ korreliert sind, tritt ein Schock in $y_2$ nie allein auf --- $y_1$ bewegt sich gleichzeitig mit.

\schritt{Schritt 2 --- Die übliche Lösung}
Man ordnet die Reihen und nimmt an, die erste reagiere \kursiv{innerhalb der Periode} nicht auf die zweite. Dann lässt sich die Korrelation eindeutig zerlegen (Cholesky-Zerlegung, Band 2).

\schritt{Schritt 3 --- Die Konsequenz}
\begin{center}
{\small
\renewcommand{\arraystretch}{1.5}
\begin{tabular}{@{}P{3.0cm}P{5.6cm}@{}}
\toprule
\textbf{Reihenfolge $(y_1,y_2)$} & Schock in $y_2$ trifft $y_1$ erst nächste Periode.\\
\textbf{Reihenfolge $(y_2,y_1)$} & Schock in $y_2$ trifft $y_1$ sofort mit $0{,}46\sigma$.\\
\bottomrule
\end{tabular}}
\end{center}

\begin{center}
\begin{minipage}{0.9\textwidth}
\textbf{Dieselben Daten, zwei Impulsantworten --- und die Reihenfolge ist eine Annahme über die Welt.}
\end{minipage}
\end{center}

\begin{center}
\begin{minipage}{0.88\textwidth}
\kursiv{\textbf{Damit ist man wieder bei Kapitel~2}: Sims wollte die Ausschlüsse der Strukturmodelle loswerden und musste sie als Reihenfolge wieder einführen. Bei Geldpolitik ist die Annahme oft vertretbar --- die Zentralbank sieht die Produktion des laufenden Quartals noch nicht. Bei Arzneipreisen und Ausgaben ist sie es weniger. Die Vektorautoregression ist ein Vorhersagewerkzeug; als Kausalwerkzeug braucht sie, was alle brauchen: eine Identifikationsannahme, die man aussprechen muss.}
\end{minipage}
\end{center}
\end{aufgabe}

\begin{geschichte}
\boxtitelg{Trunkenbold und Hund}
\Person{Clive Granger} kam auf die Kointegration, als er 1981 zeigen wollte, dass die Fehlerkorrekturmodelle von \Person{David Hendry} und \Person{Denis Sargan} mathematisch unhaltbar seien --- und beim Versuch das Gegenteil bewies.

\begin{center}
\begin{minipage}{0.9\textwidth}
\kursiv{Das Bild, das sich durchsetzte, stammt von \Person{Michael Murray}: Ein Betrunkener und sein Hund, beide ohne Ziel, beide Irrfahrten --- \textbf{aber der Hund kommt zurück, wenn der Herr ruft, und der Abstand zwischen ihnen bleibt beschränkt.} Wer wissen will, wo der Hund ist, suche beim Herrn. Der Aufsatz von \Person{Engle} und \Person{Granger} erschien 1987 und ist einer der meistzitierten der Ökonometrie; der Nobelpreis kam 2003.}
\end{minipage}
\end{center}

\schritt{Sims und die VAR}
\begin{center}
\begin{minipage}{0.9\textwidth}
\kursiv{\Person{Christopher Sims} veröffentlichte 1980 \kursiv{Macroeconomics and Reality} --- die Kritik aus Kapitel~2 und den Vorschlag der Vektorautoregression. \textbf{Er erhielt 2011 den Nobelpreis gemeinsam mit \Person{Thomas Sargent}}, dessen Kritik aus der anderen Richtung kam: dass Modelle, die Erwartungen ignorieren, für Politikbewertung untauglich sind. Kapitel~13 nimmt beide auf.}
\end{minipage}
\end{center}

\schritt{Johansen}
\begin{center}
\begin{minipage}{0.9\textwidth}
\kursiv{\Person{Søren Johansen} gab 1988 das Verfahren für mehrere Reihen an --- aus Kopenhagen, wo achtzig Jahre zuvor Erlang die Warteschlangen begründet hatte. \textbf{Sein Rangtest ist die Standardmethode für Systeme von drei und mehr Reihen} und ein seltenes Beispiel dafür, dass lineare Algebra aus Band 2 eine ökonometrische Frage unmittelbar beantwortet: Der Rang einer Matrix ist die Zahl der Gleichgewichte.}
\end{minipage}
\end{center}
\end{geschichte}

\begin{forschung}
\boxtitelp{Kointegration und VAR heute}
\begin{center}
{\small
\renewcommand{\arraystretch}{1.45}
\begin{tabular}{@{}P{3.2cm}P{5.0cm}@{}}
\toprule
\textbf{Identifikation von VARs} & Vorzeichenrestriktionen, externe Instrumente, erzählte Schocks statt Reihenfolgen.\\
\textbf{Große VARs} & Hunderte Reihen mit Schrumpfung --- Band 22.\\
\textbf{Lokale Projektionen} & Impulsantworten direkt regressieren, ohne VAR.\\
\bottomrule
\end{tabular}}
\end{center}

\schritt{Zur ersten Zeile}
\begin{center}
\begin{minipage}{0.88\textwidth}
\kursiv{Statt einer Reihenfolge nimmt man heute oft ein externes Instrument für den Schock --- \textbf{etwa die Überraschung in den Zinserwartungen während der dreißig Minuten um eine Zentralbankentscheidung.} Das ist Kapitel~3 in Zeitreihenform, und es verbindet die beiden Hälften dieses Bandes: Die Identifikationsfrage der Querschnittsökonometrie und die Dynamik der Zeitreihen treffen sich in der Frage, was ein Schock ist.}
\end{minipage}
\end{center}
\end{forschung}

\begin{merksatz}
\boxtitel{Was Sie jetzt können}
\begin{center}
\renewcommand{\arraystretch}{1.5}
\begin{tabular}{@{}P{3.0cm}P{5.6cm}@{}}
\toprule
\textbf{Definition 12.1} & stationärer Gleichgewichtsfehler\\
\textbf{Satz 12.2} & Superkonsistenz mit Rate $1/T$\\
\textbf{Satz 12.3} & Residuen testen, strengere Grenze\\
\textbf{Satz 12.4} & kurz, lang und Anpassung getrennt\\
\textbf{VAR} & Vorhersage ohne Ausschlüsse, Kausalität mit\\
\bottomrule
\end{tabular}
\end{center}

\begin{center}
\begin{minipage}{0.88\textwidth}
\kursiv{Damit endet Teil~IV. Die Zeitreihenmodelle sagen gut vorher --- besser als die Strukturmodelle, wie Kapitel~10 andeutete. Teil~V fragt, ob gute Vorhersage genügt, um eine Politik zu bewerten, und findet einen Grund, warum sie es nicht tut.}
\end{minipage}
\end{center}
\end{merksatz}
